% TERMINA ESA COSA
\pagebreak
\section{Justificación}

A esto se le suma el hecho de que pese a lo que \textcite{Herat2020} detalla
sobre las soluciones proporcionadas por las empresas, \textcite{Hartl2023}
encontraron que no todas las personas son conscientes de este hecho, y por ende
sus productos se añaden a los valores de los desperdicios eléctricos y
electrónicos. Considerando que en la mayoría de los casos se asegura un periodo
de tiempo para mantener a flote la economía de las firmas, al ver que los
usuarios tienen el deseo de comprar un aparato electrónico de mayor calidad una
vez se presentan fallos. Sin embargo, no todos se encuentran con la posibilidad
de reemplazar una máquina en periodos de tiempo, en particular si estos son
cortos, cosa que resulta costosa para los usuarios.

Por lo tanto, la idea de la aplicación a desarrollar se basa principalmente en
reducir los riesgos que se generan, y según lo expresa \textcite{Mosquera2019},
al realizar varias de las actividades que involucran directamente el
desconocimiento de la información que se suele presentar por parte los nuevos
usuarios, muchos atacantes aprovechan los vacíos de seguridad que los presenta
y por ende, hacerlos quedar expuestos a varias amenazas que atentan con la
privacidad, la identidad, el estado de una organización o las infraestructuras
críticas. \parencite[e.g.,][]{CanoMartinez2022} 
